\section{Introduction}
Oscillating heat pipes (OHPs) achieve high effective conductivity through self-excited two-phase oscillation, but the
same mechanism makes their behaviour difficult to model from first principles: the flow regime, the distribution of
liquid slugs and vapour plugs, and the evaporation--condensation dynamics are only partially observable through wall
temperatures. Working-fluid comparison is therefore usually experimental: one apparatus per fluid, per fill ratio, per
geometry. Data-driven surrogates trained on such measurements interpolate well but are known to fail when asked to
extrapolate to unseen operating conditions -- precisely the question a designer asks (``what if we run hotter?'').

Grey-box identification offers a middle road: impose the physically mandatory structure (energy conservation on a
small number of lumped nodes) and learn the few effective parameters from data. Physics-informed neural networks
(PINNs) provide a convenient machinery for this, embedding the ODE residual in the loss through automatic
differentiation and optimising physical parameters jointly with a trial solution. What PINNs add over classical
shooting-based parameter fitting for such small, well-posed ODEs is, however, an open and often glossed-over question.

This paper contributes a deliberately honest case study on a public dataset of a helically coiled OHP (HCOHP)
\cite{dataset,dib} operated with three working fluids:
\begin{enumerate}
  \item a two-node grey-box model, vessel-temperature-driven, identified by PINN and by ODE shooting, validated
        leave-one-run-out with extrapolation both up and down in heating level, against IC-fair baselines:
        trivial, linear, and nonlinear black boxes (MLP, GP, SINDy, autoregressive GRU);
  \item evidence that a correctly structured linear model, not the PINN optimiser or the physical
        parameterisation, carries the accuracy: shooting, one-node, and even an unconstrained per-fluid linear
        state-space extrapolate at \SIrange{0.3}{0.9}{\kelvin} while nonlinear black boxes degrade
        (GRU \SI{5}{\kelvin}, others \SIrange{24}{87}{\kelvin}; trivial gap baseline \SI{3}{\kelvin});
  \item a dimensionless coupling ratio $R^*=a_f/\kappa_f$ that is invariant to the sampling interval,
        seed- and method-stable, ranks the fluids identically in every fit, and is recoverable from plateau
        ratios alone -- with the ethanol drift consistent with the dry-out reported in the source data;
  \item negative results reported as such: a temperature-dependent pipe conductance $\kappa(T_e)$ is
        identifiable within each run (an initial ``not identifiable'' verdict was an optimiser artefact,
        corrected) but is a heating-level effect that does not transfer across runs, so the deployed model
        keeps $\kappa$ constant; per-fluid condenser loss rates do not pay for their parameters;
  \item a consistency check on an untrained channel: adiabatic-section temperatures are reproduced as a static
        mixture of the two identified nodes to \SI{<0.5}{\kelvin} on held-out EOHP/WOHP runs (MOHP is harder);
  \item an exported ONNX simulator used for model-implied scenario illustrations (tracking vs.\ ripple rejection
        along $R^*$) with seed-ensemble bands -- explicitly not fluid recommendations, since the model carries no
        heat-flux or dry-out limits.
\end{enumerate}
